% This is text associated with the open textbook "Open Electromagnetics"
% See immediately below for author, date
% This work is licensed under the Creative Commons Attribution-ShareAlike 4.0 International License. (CC BY-SA 4.0) 
% To view a copy of the license, visit http://creativecommons.org/licenses/by-sa/4.0/

%ORB TITLE Comparison of Static and Time-Varying Electromagnetics 
%ORB AUTHOR 0        # S.W. Ellingson, Virginia Tech (ellingson@vt.edu)
%ORB DATE 20180807   # yyyymmdd
%ORB ABSTRACT Comparison of Static and Time-Varying Electromagnetics 

\label{m0013_Comparison_of_Static_and_Time-Varying_Electromagnetics} % <-- Must match name of this file and module directory name.  Don't mess with this!

Students encountering time-varying electromagnetic fields for the first time have usually been exposed to electrostatics and magnetostatics already.
These disciplines exhibit many similarities as summarized in Table~\ref{m0013_tComparison}.
The principles of time-varying electromagnetics presented in this table are all formally introduced in other sections; the sole purpose of this table is to point out the differences. 
We can summarize the differences as follows:
%
%ORB HIGHLIGHT
\begin{mdframed}[backgroundcolor=yellow!20,nobreak=true] 
Maxwell's Equations in the general (time-varying) case include extra terms that do not appear in the equations describing electrostatics and magnetostatics.
These terms involve time derivatives of fields and describe coupling between electric and magnetic fields.
\end{mdframed}  
%
%\onecolumn
\begin{table*}[t]
\begin{center}
\begin{tabular}{lll}
~ & {\bf Electrostatics /} & {\bf Time-Varying} \\
~ & {\bf Magnetostatics}   & {\bf (Dynamic)} \\
\hline\hline
Electric \& magnetic & independent  & {\color{blue}possibly coupled} \\
fields are...       & ~  & ~ \\
\hline
Maxwell's Eqns. & $\oint_{\mathcal S}{{\bf D}\cdot d{\bf s}} = Q_{encl}$  & $\oint_{\mathcal S}{{\bf D}\cdot d{\bf s}} = Q_{encl}$ \\
~~(integral)   & $\oint_{\mathcal C}{{\bf E}\cdot d{\bf l}} = 0$         & $\oint_{\mathcal C}{{\bf E}\cdot d{\bf l}} = {\color{blue}-\frac{\partial}{\partial t}\int_{\mathcal S}{\bf B}\cdot d{\bf s}}$ \\
               & $\oint_{\mathcal S}{{\bf B}\cdot d{\bf s}} = 0$         & $\oint_{\mathcal S}{{\bf B}\cdot d{\bf s}} = 0$ \\
               & $\oint_{\mathcal C}{{\bf H}\cdot d{\bf s}} = I_{encl}$  & $\oint_{\mathcal C}{{\bf H}\cdot d{\bf l}} = I_{encl} {\color{blue} + \int_{\mathcal S}\frac{\partial}{\partial t}{\bf D}\cdot d{\bf s}}$ \\
\hline
Maxwell's Eqns.     & $\nabla\cdot{\bf D}=\rho_v$   & $\nabla\cdot{\bf D}=\rho_v$ \\
~~(differential)   & $\nabla\times{\bf E}=0$       & $\nabla\times{\bf E}={\color{blue}-\frac{\partial}{\partial t}{\bf B}}$ \\
                   & $\nabla\cdot{\bf B}=0$        & $\nabla\cdot{\bf B}=0$ \\
                   & $\nabla\times{\bf H}={\bf J}$ & $\nabla\times{\bf H}={\bf J} {\color{blue}+ \frac{\partial}{\partial t}{\bf D}}$ \\
\hline
\end{tabular}
\end{center}
\label{m0013_tComparison}
\caption{Comparison of principles governing static and time-varying electromagnetic fields.  Differences in the time-varying case relative to the static case are highlighted in {\color{blue}blue}.
\index{Maxwell's equations} 
}
\end{table*}
%\twocolumn

The coupling between electric and magnetic fields in the time-varying case has one profound consequence in particular. 
It becomes possible for fields to continue to exist even after their sources -- i.e., charges and currents -- are turned off.  
What kind of field can continue to exist in the absence of a source?  
Such a field is commonly called a \emph{wave}.
Examples of waves include signals in transmission lines and signals propagating away from an antenna. 


%-------------------------------------------------------

{\bf Additional Reading:}
\begin{itemize}
\item ``\href{https://en.wikipedia.org/wiki/Maxwell's_equations}{Maxwell's Equations}'' on Wikipedia.
\end{itemize}
